% AUTHORED
\section{Math \& tooling prerequisites}\label{sec:00}
\noindent\textbf{Setup module, no quiz shard.}

This course uses a small amount of linear algebra, calculus, probability, and
Python tooling over and over. The goal is not to make the notation impressive;
the goal is to keep every model component easy to inspect. When a later
milestone says a tensor has shape $B \times T \times C$, you should know what
each letter means. $B$ is the batch size. $T$ is the number of time steps. $C$
is the channel (or feature) dimension carried at each token position.

The course builds GPT (Generative Pre-trained Transformer) in four passes.
First, you learn the notation and
training mechanics that every later module reuses. Then you build trainable
neural-network primitives: linear layers, nonlinearities, losses, gradients, and
optimizers. After that, you add sequence primitives: token embeddings,
self-attention, causal masking, multiple heads, and position information.
Finally, you assemble those pieces into a decoder-only language model that can
train on next-token prediction and sample new text.

\subsection*{Vectors, matrices, and shapes}
A vector is an ordered list of numbers. A matrix is a rectangular table of
numbers. A linear layer usually computes
\[
  Y = XW + b,
\]
where $X \in \mathbb{R}^{B \times d_{\text{in}}}$ contains a batch of examples,
$W \in \mathbb{R}^{d_{\text{in}} \times d_{\text{out}}}$ contains trainable
weights, $b \in \mathbb{R}^{d_{\text{out}}}$ is broadcast across the batch, and
$Y \in \mathbb{R}^{B \times d_{\text{out}}}$ is the result. Most bugs in this
course are shape bugs, so write shapes next to new formulas before you write
code.

\subsection*{Derivatives and the chain rule}
Training means changing parameters to reduce a scalar loss $L$. A derivative
tells you how a small change in one number changes another number. For a
composition $L=f(g(x))$, the chain rule says
\[
  \frac{dL}{dx} = \frac{dL}{dg}\frac{dg}{dx}.
\]
Backpropagation is the chain rule applied systematically to every operation in
a computation graph. You will first do this by hand for simple networks, then
encode it in the autograd milestone.

\subsection*{Probability and log loss}
A language model predicts a probability distribution over the next token. A
distribution has nonnegative entries that sum to one. The loss used throughout
the course is negative log-likelihood: if the correct next token has probability
$p$, the loss is $-\log p$. This heavily penalizes confident wrong predictions
and rewards placing probability mass on the observed token.

\subsection*{Numerical habits}
Floating-point numbers have limits. Large exponentials overflow; tiny
probabilities can underflow; subtracting nearly equal values can lose precision.
Three later techniques follow the same idea: the stable softmax (concept 03),
gradient clipping (optimization), and careful masking (attention). In each case
the goal is the same: keep the mathematical result correct while computing it in
a range the machine can represent.

\subsection*{Tooling loop}
Use the repository the way you would use a lab notebook. Read the module, run
the small exercise, take the quiz, then advance the project tests. The concept
exercises are small, self-contained numpy tasks. The project track is cumulative and lives under
\texttt{project/dlgpt}. The fastest way to learn is to run tests frequently and
inspect intermediate shapes rather than writing a large block of code at once.
Concept 00 is setup-only, concepts 01--13 have quiz shards, and P0--P9 are the
cumulative project milestones.

\subsection*{Project milestones (P0--P9)}
Later modules refer to numbered project milestones. They build on one another in
order, each reusing the concept named in parentheses:
\begin{itemize}
  \item \textbf{P0} -- Bootstrap: character tokenizer and batch loader (provided, already written).
  \item \textbf{P1} -- Autograd engine (concept 05).
  \item \textbf{P2} -- Neural-network primitives: linear layer, activation, loss, optimizers (concepts 03, 04, 06).
  \item \textbf{P3} -- Embeddings and a character language-model baseline (concept 07).
  \item \textbf{P4} -- Single-head self-attention (concept 08).
  \item \textbf{P5} -- Causal masking (concept 09).
  \item \textbf{P6} -- Multi-head attention (concept 10).
  \item \textbf{P7} -- Positional encoding (concept 11).
  \item \textbf{P8} -- Transformer block (concept 12).
  \item \textbf{P9} -- Decoder-only GPT and sampling (concept 13).
\end{itemize}
Only P0 comes already written; you build P1--P9 by making their failing tests pass.

\medskip
\noindent\textit{Go deeper:} Goodfellow, Bengio \& Courville, \emph{Deep
Learning}, Ch.~2--4 (\texttt{https://www.deeplearningbook.org}) covers the
linear-algebra, probability, and numerical-computation background this course
relies on.
